\section{Erd\H{o}s--Hajnal reciprocal cycle-length sum (Problem 65) --- Round 12}

\subsection*{1) ROUND-12 OBJECTIVE}

\textbf{Path chosen: (C) obstruction/correction.}
Round~11 corrected the literature picture by importing a large-$d$ threshold theorem, but it still treated the second Erd\H{o}s--Hajnal question as though the main remaining task were to prove a universal exact-edge minimiser statement.

In this round I resolve the statement \emph{as written} by showing that it is false (more precisely, fatally underspecified) under the natural exact-edge interpretation used throughout the earlier rounds. The obstruction is simple and rigorous: for infinitely many pairs $(n,m)$ there is \emph{no} complete bipartite graph on $n$ vertices with exactly $m$ edges. Hence the literal statement
\begin{quote}
``for fixed $(n,m)$, the quantity $S(G)$ is minimised when $G$ is a complete bipartite graph''
\end{quote}
cannot be true as a universal theorem over exact edge classes.

I then state the minimal corrected exact-edge version and record the currently verified regimes for that corrected problem.

\subsection*{2) Round-11 FOUNDATION USED}

I use the following vetted Round~11 inputs.
\begin{itemize}
\item[(F1)] For integers $a,b\ge 1$, the complete bipartite graph $K_{a,b}$ has
\[
\mathcal C(K_{a,b})=\{4,6,\dots,2\min\{a,b\}\},
\qquad
S(K_{a,b})=\frac12\bigl(H_{\min\{a,b\}}-1\bigr).
\]
\item[(F2)] Round~10 proved the exact line $m=2n-4$ for $7\le n\le 51$:
\[
\min\{S(G): |V(G)|=n,\ |E(G)|=2n-4\}=\frac14,
\]
attained by $K_{2,n-2}$.
\item[(F3)] Round~11 imported the survey theorem that there exists $d_0$ such that for every $d\ge d_0$, among $n$-vertex graphs with at least $d(n-d)$ edges, $S(G)$ is minimised exactly by $K_{d,n-d}$.
\end{itemize}

\subsection*{3) NEW INSIGHT / TOOL (ROUND-12)}

The new observation is a basic arithmetic obstruction.
For a complete bipartite graph on $n$ vertices, the edge count must be of the form
\[
ab
\qquad\text{with }a+b=n,
\]
that is,
\[
a(n-a)\qquad(0\le a\le n).
\]
These values do \emph{not} cover all integers between $0$ and $\lfloor n^2/4\rfloor$.
In particular, when $n=2t+1$ is odd, the largest possible complete-bipartite edge count is
\[
t(t+1),
\]
while the next largest is
\[
(t-1)(t+2)=t(t+1)-2.
\]
So the value
\[
m=t(t+1)-1
\]
is \emph{not} realised by any complete bipartite graph on $2t+1$ vertices.
This yields an infinite counterexample family to the literal exact-edge formulation.

\subsection*{4) ATTACK PLAN (ROUND-12)}

The exact issue left vague in earlier rounds is the meaning of the sentence
\begin{quote}
``Is the sum $\sum 1/a_i$ minimised when $G$ is a complete bipartite graph?''
\end{quote}
from Problem~65.

I now interpret it in the same exact-edge sense used throughout Rounds~2--11:
for fixed integers $(n,m)$, minimise $S(G)$ over all graphs with $|V(G)|=n$ and $|E(G)|=m$.
Under that interpretation, the plan is:
\begin{enumerate}
\item prove that infinitely many exact edge counts $m$ are not achievable by any complete bipartite graph on $n$ vertices;
\item choose such a pair $(n,m)$ with the feasible class nonempty;
\item conclude that the minimum over that class is attained by some graph, but by no complete bipartite graph, disproving the literal statement;
\item formulate the minimal corrected statement: restrict to pairs $(n,m)$ with $m=a(n-a)$.
\end{enumerate}

\subsection*{5) WORK (ROUND-12)}

As before, for a graph $G$ let
\[
\mathcal C(G):=\{\ell\ge 3:\text{$G$ contains a cycle of length $\ell$}\},
\qquad
S(G):=\sum_{\ell\in\mathcal C(G)}\frac1\ell.
\]

\subsubsection*{5.1. Complete bipartite edge counts on a fixed vertex set}

\textbf{Lemma 5.1.}
A complete bipartite graph on $n$ vertices has exactly $a(n-a)$ edges for some integer $a$ with $0\le a\le n$.
Conversely, for each such $a$, the graph $K_{a,n-a}$ has exactly $a(n-a)$ edges.

\medskip
\textit{Proof.}
A complete bipartite graph on $n$ vertices has a bipartition into parts of sizes $a$ and $n-a$, and every possible cross-edge is present, so the edge count is exactly $a(n-a)$.
The converse is immediate from the definition of $K_{a,n-a}$.
\hfill$\square$

\medskip
\textbf{Lemma 5.2.}
Let $n=2t+1$ with $t\ge 1$.
Among complete bipartite graphs on $n$ vertices, the largest edge count is $t(t+1)$, attained by $K_{t,t+1}$ and $K_{t+1,t}$, and the next largest edge count is $t(t+1)-2$, attained by $K_{t-1,t+2}$ and $K_{t+2,t-1}$.
In particular, no complete bipartite graph on $2t+1$ vertices has exactly
\[
t(t+1)-1
\]
edges.

\medskip
\textit{Proof.}
By Lemma~5.1 the possible edge counts are
\[
f(a):=a(2t+1-a)\qquad(0\le a\le 2t+1).
\]
This is a concave quadratic in $a$, symmetric about $a=t+\tfrac12$.
Therefore its maximum over integers is attained at $a=t$ and $a=t+1$, where
\[
f(t)=f(t+1)=t(t+1).
\]
The next closest integers are $a=t-1$ and $a=t+2$, and there
\[
f(t-1)=(t-1)(t+2)=t(t+1)-2.
\]
By concavity, all remaining values are at most $t(t+1)-2$.
Hence $t(t+1)-1$ is not realised.
\hfill$\square$

\subsubsection*{5.2. Infinite counterexample family to the literal exact-edge statement}

\textbf{Theorem 5.3.}
Under the exact-edge interpretation of Problem~65, the statement
\begin{quote}
``for fixed $(n,m)$, the minimum of $S(G)$ is attained when $G$ is complete bipartite''
\end{quote}
is false.
In fact there is an infinite family of counterexamples: for every integer $t\ge 2$, let
\[
n=2t+1,
\qquad
m=t(t+1)-1.
\]
Then among graphs with $n$ vertices and $m$ edges, the minimum of $S(G)$ is attained by some graph, but by no complete bipartite graph.

\medskip
\textit{Proof.}
Fix $t\ge 2$ and let $(n,m)=(2t+1,\,t(t+1)-1)$.
First, the feasible class
\[
\mathcal G_{n,m}:=\{G:|V(G)|=n,\ |E(G)|=m\}
\]
is nonempty, because one may simply choose any set of $m$ edges from the complete graph $K_n$.
Since $\mathcal G_{n,m}$ is finite, the set of values
\[
\{S(G):G\in\mathcal G_{n,m}\}
\]
has a minimum, attained by at least one graph $G_0\in\mathcal G_{n,m}$.

By Lemma~5.2, no complete bipartite graph on $n=2t+1$ vertices has exactly $m=t(t+1)-1$ edges.
Therefore no member of $\mathcal G_{n,m}$ is complete bipartite.
In particular, the minimiser $G_0$ is not complete bipartite.
So the literal exact-edge statement fails for $(n,m)$.
Because this construction works for every $t\ge 2$, it gives infinitely many counterexamples.
\hfill$\square$

\subsubsection*{5.3. Minimal corrected statement}

The failure above is not a subtle extremal phenomenon; it is a missing-hypothesis problem.
The minimal corrected exact-edge formulation is therefore:

\begin{quote}
\textbf{Corrected exact-edge question.}
For integers $n,m$ such that there exists an integer $a$ with
\[
0\le a\le n
\qquad\text{and}\qquad
m=a(n-a),
\]
is it true that
\[
\min\{S(G):|V(G)|=n,\ |E(G)|=m\}=S(K_{a,n-a}),
\]
and that the minimum is attained by $K_{a,n-a}$?
\end{quote}

This is the strongest exact-edge correction that changes only the missing existence hypothesis.

There is also a natural threshold version:
\begin{quote}
\textbf{Threshold question.}
For fixed $n$ and $d$, among all $n$-vertex graphs with at least $d(n-d)$ edges, is $S(G)$ minimised by $K_{d,n-d}$?
\end{quote}

By the checked 2026 survey, this threshold question is already solved for all sufficiently large $d$, while the exact-edge corrected question is solved in at least the following regimes:
\begin{itemize}
\item $a=2$ and $7\le n\le 51$ by Round~10;
\item all sufficiently large $a$ via the large-$d$ threshold theorem imported in Round~11.
\end{itemize}

\subsection*{6) ADVERSARIAL VERIFICATION}

\begin{itemize}
\item \textbf{Does Theorem~5.3 really disprove the literal statement?}
Yes. The feasible class is finite and nonempty, so a minimiser exists. If the theorem's conclusion were true for all $(n,m)$ in the exact-edge sense, that minimiser would have to be complete bipartite. But for the displayed infinite family no complete bipartite graph exists in the class at all.

\item \textbf{Did we need to know the actual minimum value?}
No. Nonexistence of a complete bipartite graph in the exact class is enough to refute the literal attainment claim.

\item \textbf{Could the original wording have meant something else?}
Yes. That is exactly why this is a path-(C) correction: the statement is fatally ambiguous unless one specifies either the exact-edge class or a threshold class. The corrected statements in \S5.3 remove that ambiguity.

\item \textbf{Why is the corrected exact-edge formulation minimal?}
Because the only added hypothesis is the necessary existence condition $m=a(n-a)$ ensuring that a complete bipartite graph with the prescribed $(n,m)$ parameters actually exists.
\end{itemize}

\subsection*{7) FINAL (EXACTLY ONE)}

\textbf{FULL SOLUTION --- COUNTEREXAMPLE / DISPROOF.}

The second statement of Problem~65 is false as written under the natural exact-edge interpretation used in the previous rounds.
An infinite counterexample family is
\[
(n,m)=(2t+1,\,t(t+1)-1)\qquad(t\ge 2),
\]
for which the optimisation class is nonempty but contains no complete bipartite graph at all.
Hence the claim that the minimum is attained ``when $G$ is complete bipartite'' cannot hold universally.

The minimal corrected exact-edge statement is:
\begin{quote}
if $m=a(n-a)$ for some integer $a$, is the minimum of $S(G)$ over all $n$-vertex $m$-edge graphs attained by $K_{a,n-a}$?
\end{quote}
This corrected question remains only partially solved, with proven regimes including $a=2$ for $7\le n\le 51$ and all sufficiently large $a$ in the threshold formulation.

\subsection*{8) COMPLETION ESTIMATE (MANDATORY)}

\textbf{COMPLETION: 100\%}

\subsection*{9) REFERENCES}

\begin{thebibliography}{9}

\bibitem{Problem65Page}
T.~F.~Bloom,
\emph{Erd\H{o}s Problem \#65}, online problem page.
Checked 2026-03-08.
The page states the problem as
``Is the sum $\sum 1/a_i$ minimised when $G$ is a complete bipartite graph?''
and notes the 2026 survey update.

\bibitem{Montgomery2025}
R.~Montgomery,
\emph{Cycles and expansion in graphs},
EMS Magazine \textbf{138} (2025/2026 issue), 5--12.
Used here only for the checked large-$d$ threshold theorem mentioned in \S5.3.

\end{thebibliography}
